\section{Introdução}

A formulação de políticas públicas municipais no Brasil ocorre em um ambiente de forte heterogeneidade institucional, textual e tecnológica. Municípios com problemas semelhantes produzem proposições com vocabulário distinto, graus diferentes de detalhamento e portais legislativos com variados níveis de padronização. Isso dificulta a reutilização de experiências anteriores, reduz a comparabilidade entre iniciativas e limita o uso sistemático de evidências legislativas no processo decisório.

Este trabalho estuda se é possível transformar o histórico de Projetos de Lei municipais em um acervo pesquisável de ações públicas orientadas a problema, integrando textos legislativos, recuperação semântica e indicadores oficiais. Em uma frase, a pergunta central de pesquisa deste artigo é: \textit{é possível construir um pipeline reproduzível para descoberta, organização e recomendação de políticas públicas municipais a partir de PLs históricos e indicadores oficiais?}

O trabalho adota um posicionamento híbrido, na interseção entre governo digital municipal, mineração de texto e recuperação de informação aplicada. O objetivo principal é propor e validar um método operacionalizado em uma plataforma web para gestores e analistas não técnicos.

As contribuições do trabalho são:
\begin{itemize}[leftmargin=*,nosep]
  \item \textbf{C1}: método reproduzível para descoberta de instâncias do SAPL e extração automatizada de PLs em escala nacional;
  \item \textbf{C2}: transformação de ementas jurídicas em ações operacionais para reduzir ruído linguístico na recuperação;
  \item \textbf{C3}: busca semântica de ações a partir de consultas em linguagem natural;
  \item \textbf{C4}: associação exploratória entre grupos de PLs e indicadores municipais oficiais ao longo do tempo;
  \item \textbf{C5}: plataforma web que integra as etapas anteriores em uma interface única.
\end{itemize}

Na implementação atual, todo o pipeline já foi construído, mas as etapas têm níveis diferentes de validação empírica. Coleta e tradução de ementas estão mais maduras; recuperação semântica, agrupamento e análise com indicadores ainda têm avaliação mais exploratória.
